
\subsection{与乘法结构的相容性}
\label{subsec:6-4-5}

现在证明定理~\ref{thm:D} 的 \textup{(D.4)}，
从而完成该定理的证明。
在第 3.7.7 小节定义了弗洛尔上同调上的乘积结构
\[
\mathfrak m_2:
HF((L,\mathbf b_2),(L,\mathbf b_1))
\otimes HF((L,\mathbf b_1),(L,\mathbf b_0))
\longrightarrow HF((L,\mathbf b_2),(L,\mathbf b_0)).
\]
它由算子
$\mathfrak m_2^{\mathbf b_2,\mathbf b_1,\mathbf b_0}
:\widehat C(L)\otimes\widehat C(L)\to\widehat C(L)$
诱导；其中
\begin{equation}
\mathfrak m_2^{\mathbf b_2,\mathbf b_1,\mathbf b_0}(x,y)
=\mathfrak m^{\boldsymbol b}
\bigl(e^{b_2},x,e^{b_1},y,e^{b_0}\bigr),
\tag{6.4.27}\label{eq:6-4-27}
\end{equation}
并且
$\mathbf b_i=(\boldsymbol b,b_i)$
（见 (3.7.90)）。
由 (6.4.27) 易见
\begin{equation}
\mathfrak m_2^{\mathbf b_2,\mathbf b_1,\mathbf b_0}
\bigl(\mathfrak F^{q_1}\widehat C(L)
\otimes\mathfrak F^{q_2}\widehat C(L)\bigr)
\subseteq\mathfrak F^{q_1+q_2}\widehat C(L).
\tag{6.4.28}\label{eq:6-4-28}
\end{equation}
因此，$\mathfrak m_2$
保持 $HF((L,\mathbf b_j),(L,\mathbf b_i))$
上的滤过 $\mathfrak F$。

为在 $E_r$ 上定义乘积，
只需证明下列引理。

\begin{lemma}{6.4.29}\label{lem:6-4-29}
\begin{align}
\mathfrak m_2^{\mathbf b_2,\mathbf b_1,\mathbf b_0}
\bigl(Z_r^{p,q}\otimes Z_r^{p',q'}\bigr)
&\subseteq Z_r^{p+p',q+q'},
\tag{6.4.30.1}\label{eq:6-4-30-1}\\
\mathfrak m_2^{\mathbf b_2,\mathbf b_1,\mathbf b_0}
\bigl(B_r^{p,q}\otimes Z_r^{p',q'}\bigr)
&\subseteq B_r^{p+p',q+q'},
\tag{6.4.30.2}\label{eq:6-4-30-2}\\
\mathfrak m_2^{\mathbf b_2,\mathbf b_1,\mathbf b_0}
\bigl(Z_r^{p,q}\otimes B_r^{p',q'}\bigr)
&\subseteq B_r^{p+p',q+q'}.
\tag{6.4.30.3}\label{eq:6-4-30-3}
\end{align}
在 (6.4.30.1) 中，
三个 $Z_r$ 分别由
$\delta_{\mathbf b_2,\mathbf b_1}$、
$\delta_{\mathbf b_1,\mathbf b_0}$、
$\delta_{\mathbf b_2,\mathbf b_0}$ 定义；
(6.4.30.2)、(6.4.30.3)
中的 $B_r$ 亦然。
\end{lemma}

\begin{proof}
只证明 (6.4.30.2)，
其余两式同理。
回忆
\begin{equation}
\begin{aligned}
&\delta_{\mathbf b_2,\mathbf b_0}
\bigl(\mathfrak m_2^{\mathbf b_2,\mathbf b_1,\mathbf b_0}(x,y)\bigr)\\
&\quad=-\mathfrak m_2^{\mathbf b_2,\mathbf b_1,\mathbf b_0}
\bigl(\delta_{\mathbf b_2,\mathbf b_1}(x),y\bigr)
+(-1)^{\deg x}
\mathfrak m_2^{\mathbf b_2,\mathbf b_1,\mathbf b_0}
\bigl(x,\delta_{\mathbf b_1,\mathbf b_0}(y)\bigr).
\end{aligned}
\tag{6.4.31}\label{eq:6-4-31}
\end{equation}
见 (3.7.91)。
取
$x=\delta_{\mathbf b_2,\mathbf b_1}x_1+x_2
\in B_r^{p,q}$，
其中
\[
\begin{gathered}
x_1\in\mathfrak F^{q-r+2}\widehat C^{p-1},
\qquad
\delta_{\mathbf b_2,\mathbf b_1}x_1
\in\mathfrak F^q\widehat C^p,\\
x_2\in\mathfrak F^{q+1}\widehat C^p.
\end{gathered}
\]
再取 $y=y_1+y_2\in Z_r^{p',q'}$，
其中
\[
\begin{gathered}
\delta_{\mathbf b_1,\mathbf b_0}y_1
\in\mathfrak F^{q'+r-1}\widehat C^{p'+1},
\qquad
y_1\in\mathfrak F^{q'}\widehat C^{p'},\\
y_2\in\mathfrak F^{q'+1}\widehat C^{p'}.
\end{gathered}
\]
这些分解见 (6.3.19)。
由滤过相容性易见
\[
\begin{aligned}
\mathfrak m_2^{\mathbf b_2,\mathbf b_1,\mathbf b_0}
\bigl(\delta_{\mathbf b_2,\mathbf b_1}x_1,y_2\bigr)
&\in\mathfrak F^{q+q'+1}\widehat C^{p+p'}
\subseteq B_r^{p+p',q+q'},\\
\mathfrak m_2^{\mathbf b_2,\mathbf b_1,\mathbf b_0}(x_2,y_1)
&\in\mathfrak F^{q+q'+1}\widehat C^{p+p'}
\subseteq B_r^{p+p',q+q'},\\
\mathfrak m_2^{\mathbf b_2,\mathbf b_1,\mathbf b_0}(x_2,y_2)
&\in\mathfrak F^{q+q'+1}\widehat C^{p+p'}
\subseteq B_r^{p+p',q+q'}.
\end{aligned}
\]
另一方面，由 (6.4.31) 得\translatornote{原书实体页 392
在代入 $(x_1,y_1)$ 时
仍把此处符号指数写成 $\deg x$；
按 (6.4.31) 应取第一个输入的次数
$\deg x_1$，见疑误 A-0390。}
\[
\begin{aligned}
&\mathfrak m_2^{\mathbf b_2,\mathbf b_1,\mathbf b_0}
\bigl(\delta_{\mathbf b_2,\mathbf b_1}x_1,y_1\bigr)\\
&\quad=-\delta_{\mathbf b_2,\mathbf b_0}
\bigl(\mathfrak m_2^{\mathbf b_2,\mathbf b_1,\mathbf b_0}
(x_1,y_1)\bigr)
+(-1)^{\deg x_1}
\mathfrak m_2^{\mathbf b_2,\mathbf b_1,\mathbf b_0}
\bigl(x_1,\delta_{\mathbf b_1,\mathbf b_0}(y_1)\bigr).
\end{aligned}
\]
由 (6.4.28) 的滤过相容性，
\[
\mathfrak m_2^{\mathbf b_2,\mathbf b_1,\mathbf b_0}(x_1,y_1)
\in\mathfrak F^{q+q'-r+2}\widehat C^{p+p'-1},
\]
因此其边界满足
\[
\begin{aligned}
&\delta_{\mathbf b_2,\mathbf b_0}
\bigl(\mathfrak m_2^{\mathbf b_2,\mathbf b_1,\mathbf b_0}
(x_1,y_1)\bigr)\\
&\quad\in
\delta_{\mathbf b_2,\mathbf b_0}
\bigl(\mathfrak F^{q+q'-r+2}
\widehat C^{p+p'-1}\bigr)
\cap\mathfrak F^{q+q'}\widehat C^{p+p'}
\subseteq B_r^{p+p',q+q'}.
\end{aligned}
\]
另一方面，
$x_1$ 与 $\delta_{\mathbf b_1,\mathbf b_0}(y_1)$
的滤过次数还给出
\[
\mathfrak m_2^{\mathbf b_2,\mathbf b_1,\mathbf b_0}
\bigl(x_1,\delta_{\mathbf b_1,\mathbf b_0}(y_1)\bigr)
\in\mathfrak F^{q+q'+1}\widehat C^{p+p'}
\subseteq B_r^{p+p',q+q'}.
\]
于是 (6.4.30.2) 得证。
\end{proof}

引理~\ref{lem:6-4-29} 表明，
$\mathfrak m_2^{\mathbf b_2,\mathbf b_1,\mathbf b_0}$
在 $E_r^{p,q}$ 上诱导乘积 $\mathfrak m_2$。
又因
$\mathfrak m_2^{\mathbf b_2,\mathbf b_1,\mathbf b_0}
\equiv\overline{\mathfrak m}_2
\pmod{T^{\lambda_0}\Lambda_{0,\mathrm{nov}}}$，
可得乘积与 \textup{(D.1)} 中的同构相容。
与 \textup{(D.2)} 中同构的相容性则由定义立即得到。
由 (6.4.31) 立刻推出
\begin{equation}
\delta_r(\mathfrak m_2(x,y))
=-\mathfrak m_2(\delta_r(x),y)
+(-1)^{\deg x}\mathfrak m_2(x,\delta_r(y)).
\tag{6.4.32}\label{eq:6-4-32}
\end{equation}
定理~\ref{thm:D} 至此证毕。
\hfill$\square$

\begin{remark}{6.4.33}\label{rem:6-4-33}
通常上同调与弗洛尔上同调之间的谱序列
应保持乘积结构这一事实，
由比兰在 \textbf{[Bir04]} 中提出。
\end{remark}

\begin{statement}{问题}{6.4.34}\label{prob:6-4-34}
我们的谱序列在何种意义下与
高阶 Massey 积
$\mathfrak m_k$（$k\ge3$）相容？
\end{statement}

上述谱序列与诱导乘积 $\mathfrak m_2$
的相容式 (6.4.32)，
使第二作者吴龙根关于单调拉格朗日环面的
马斯洛夫类所证明的结果
（\textbf{[Oh96I]} 的定理 III，
即奥丹问题\liftIndex{I216}）
得到如下改进：
原结果在较为特殊的维数限制
$\dim L\le24$ 下成立。
本证明使用谱序列相容性的想法来自
塞德尔 \textbf{[Sei01]}。
列夫·布霍夫斯基（Lev Buhovsky）
在 \textbf{[Buh06]} 中独立证明了
定理~\ref{thm:6-4-35}；
他证明了单调情形下
吴龙根谱序列的乘法结构相容性\liftIndex{I370}。
他的证明与此处的证明不同；
这里滤过 $A_\infty$ 公式起关键作用。

\begin{theorem}{6.4.35}\label{thm:6-4-35}
设 $L\subset(M,\omega)$
是单调拉格朗日环面，
存在哈密顿微分同胚 $\phi$，
使 $\phi(L)\cap L=\varnothing$。
则其极小马斯洛夫数满足
\[
\Sigma_L=2.
\]
\end{theorem}

\begin{proof}
首先，$T^n$ 可定向且为 spin，
而按定义，单调性蕴含 $\Sigma_L>0$。
用反证法。
假设 $\Sigma_L\ge3$。
由于 $T^n$ 可定向，
事实上有 $\Sigma_L\ge4$，
但证明中不需要这一点。
由单调性，
按照 \textbf{[Oh96I]}
或第 2.4.4 小节相同的论证，
$L$ 无障碍。
易见 $b=0\in\mathcal M(L)$，
故未经形变的弗洛尔微分 $\delta$
满足 $\delta^2=0$。

利用单调性作次数计数，
微分可分解为
\[
\delta
=\delta_{(0)}+\delta_{(1)}
+\cdots+\delta_{(N)}
=:\delta_{(0)}+\delta',
\]
其中
$N\le\left[\frac{n+2}{\Sigma_L}\right]$，
而
\begin{equation}
\delta_{(j)}
=\overline\delta_{(j)}
\otimes T^{j\lambda_0}e^{j\Sigma_L}.
\tag{6.4.36}\label{eq:6-4-36}
\end{equation}
这里 $\overline\delta_{(j)}:E_j\to E_j$
的次数为 $1-j\Sigma_L$。
由于 $\Sigma_L\ge3$，
当 $j\ge1$ 时此次数不大于 $-2$。

考虑 $\delta$ 在
\[
E_2=H^*(T^n;\mathbb Q)
\otimes\Lambda_{0,\mathrm{nov}}
\]
上的作用。
回忆上同调环 $H^*(T^n;\mathbb Q)$
由一次上同调类
$\alpha_1,\ldots,\alpha_n
\in H^1(L;\mathbb Q)$ 生成。
次数计数表明，
每个正能量分量都满足
\[
\overline\delta_{(j)}(\alpha_i)=0
\qquad(i=1,\ldots,n;\ j\ge1),
\]
因而 $\delta_2(\alpha_i)=0$。
\translatornote{原书实体页 393
把发生元下标与能量层下标都写成 $k$，
并仅以一个 $\overline\delta_{(k)}(\alpha_k)$
表示 $\delta_2(\alpha_k)$；
这里拆开两个独立指标并逐层说明为零，
见疑误 A-0387。}
由 (6.4.32) 可得\translatornote{原书实体页 393
在式 (6.4.37) 的两个乘积项前
都漏去了负号；
由 $\deg\alpha_i=1$ 代入 (6.4.32)
应有如下符号，最终零结论不变，
见疑误 A-0386。}
\begin{equation}
\begin{aligned}
\delta_2\bigl(\mathfrak m_2(\alpha_i,\alpha_j)\bigr)
&=-\mathfrak m_2\bigl(\delta_2(\alpha_i),\alpha_j\bigr)\\
&\quad-\mathfrak m_2\bigl(\alpha_i,\delta_2(\alpha_j)\bigr)=0.
\end{aligned}
\tag{6.4.37}\label{eq:6-4-37}
\end{equation}
另一方面，按能量考虑，
\[
\mathfrak m_2(\alpha_i,\alpha_j)
\equiv\alpha_i\cup\alpha_j.
\]
于是对任意 $i,j=1,\ldots,n$，
$\delta_2(\alpha_i\cup\alpha_j)=0$。
对 $\alpha_i$ 的各种乘积归纳应用同一论证，
谱序列在 $E_2$ 页退化，故
\[
HF^*(L;\Lambda_{0,\mathrm{nov}})
\cong H^*(T^n;\mathbb Q)
\otimes\Lambda_{0,\mathrm{nov}}.
\]
特别地，它是
$\Lambda_{0,\mathrm{nov}}$ 上的自由模，
于是
\begin{equation}
HF^*(L;\Lambda_{\mathrm{nov}})
\cong H^*(T^n;\mathbb Q)
\otimes\Lambda_{\mathrm{nov}}\ne0.
\tag{6.4.38}\label{eq:6-4-38}
\end{equation}
（此前尚未使用 $L\subset M$ 可位移的假设；
它只在证明的最后一步使用。）
另一方面，
弗洛尔上同调在哈密顿形变下不变，
且存在 $\phi$ 使
$\phi(L)\cap L=\varnothing$，
故
$HF^*(L;\Lambda_{\mathrm{nov}})=\{0\}$。
这与 (6.4.38) 矛盾，
因而 $\Sigma_L\le2$。
又因 $T^n$ 可定向，
非零整数 $\Sigma_L$ 必为偶数，
故 $\Sigma_L=2$。
\end{proof}

\pagebreak
同理可证下述定理。

\begin{theorem}{6.4.35bis}\label{thm:6-4-35bis}
设 $L\subset(M,\omega)$
为单调拉格朗日子流形，
且存在哈密顿微分同胚 $\phi$，
使 $\phi(L)\cap L=\varnothing$。
再假设下列条件之一成立：
\begin{enumerate}
\item[(1)] 环 $H^*(L;\mathbb Z_2)$
由 $H^1(L;\mathbb Z_2)$ 生成；
\item[(2)] $L$ 可定向且相对 spin，
并存在交换环 $R$，
使环 $H^*(L;R)$
由 $H^1(L;R)$ 与 $H^2(L;R)$ 生成。
\end{enumerate}
则 $L$ 的极小马斯洛夫数满足
$\Sigma_L=2$。
\end{theorem}

\begin{proof}
由于 $L$ 单调，
在 \textup{(1)} 中可取
$\mathbb Z_2$ 系数，
在具备相对 spin 定向资料的
\textup{(2)} 中可取相应的交换环 $R$，
以定义弗洛尔上同调。
\translatornote{原书实体页 394
仅凭单调性便称可用“任意系数环 $R$”；
但 (1) 只需 $\mathbb Z_2$，
(2) 另有相对 spin 假设才有一般 $R$
所需的定向，见疑误 A-0388。}
情形 \textup{(1)} 的证明
与定理~\ref{thm:6-4-35} 完全相同。
\translatornote{原书实体页 394
以“证明相同”结束情形 (1)，
但定理 6.4.35 最后还靠可定向性
排除 $\Sigma_L=1$；
本定理 (1) 未假设可定向，
上述论证尚未排除此值。
保留源结论与这一证明缺口，
见疑误 A-0389。}
情形 \textup{(2)} 中，
由可定向性知
$\Sigma_L\ne2$ 时有 $\Sigma_L\ge4$。
故在 $H^*(L;R)$
由一、二次类生成这一更弱条件下，
定理~\ref{thm:6-4-35} 的证明仍然适用。
\end{proof}

\noindent\textbf{定理 6.1.4 的证明。}
论证与定理 D 的证明类似。
事实上，由第 3.7.5
与第 6.3.3 小节的结果，
可构造满足 (6.1.5.1)、(6.1.5.2)
的谱序列。
(6.1.5.3) 的证明
与上面对 \textup{(D.4)} 的证明类似。
细节留给读者。
\hfill$\square$
